\documentclass[10pt,a4paper]{amsart}
\usepackage[margin=19mm]{geometry}
\usepackage{amsmath,amssymb,mathtools}
\usepackage[scr=rsfs,cal=euler]{mathalfa}
\usepackage{xfrac}
\usepackage[unicode,hidelinks,bookmarks=false]{hyperref}
\newcommand{\ie}{i.e.\ }
\newcommand{\cf}{cf.\ }
\newcommand{\resp}{respectively\ }
\newcommand{\Subsection}{\subsection}
\providecommand{\nobreakdash}{\nobreak\hbox{-}}
\setlength{\emergencystretch}{2em}
\multlinegap0pt
\usepackage{amssymb}
\usepackage{tikz}
\usepackage{tcolorbox}
\usepackage[shortlabels]{enumitem}
\usepackage{xcolor}
\newtheorem{lem}{Lemma}
\usepackage{cleveref}
\crefname{lem}{Lemma}{Lemmas}
\renewcommand{\P}{{\mathcal P}}
\newcommand{\Q}{{\mathcal Q}}
\def\Seg{\operatorname{Seg}}
\newcommand{\uSdw}{\rotatebox[origin=c]{180}{$\Delta$}}
\title{Constructions of Macaulay Posets and Macaulay Rings}
\author{Penelope Beall, Erenay Boyali, Nancy Chen, Ellen Chlachidze, Trong Toan Dao, Frederic Garvey, Mitchell Johnson, Yu Olivier Li, Nikola Kuzmanovski, Kelvin Ma, Treanungkur Mal, Rukshan Marasinghe Mudiyanselage, Quinlan Mayo, Nava Minsky-Primus, Alexandra Seceleanu, Sriram Veerapaneni}
\date{Source: arXiv:2502.15166v2 (CC BY 4.0)}
\begin{document}
\maketitle

\noindent\emph{Source and licence.} Constructions of Macaulay Posets and Macaulay Rings. Version arXiv:2502.15166v2 (CC BY 4.0), distributed by arXiv under the Creative Commons Attribution 4.0 International licence, http://creativecommons.org/licenses/by/4.0/, as recorded in the arXiv licence metadata for this exact version and shown on the abstract page https://arxiv.org/abs/2502.15166v2. Full licence notice: \texttt{licenses/arxiv-2502.15166-CC-BY-4.0.txt}.\par\smallskip

\noindent\textit{Editorial scope.} Counted unit: the article's lem Mac (source0.tex lines 1836--1838). The article's complete original proof of it is delivered with it (source0.tex lines 1840--1871, one \texttt{proof} environment of 32 lines). No mathematics is written, repaired or re-derived: the statement and the proof are the article's own lines, in the article's own notation, and no retained line is substituted or re-flowed.  Retained uncounted as the source-local context the proof needs: lines 1818--1820.  Nothing was omitted from any retained span, so the package carries no omission record.  No retained line contains a citation, so the wrapper carries no bibliography and no bibliography entry.  No retained line contains a figure environment, an image-inclusion command or drawing code, so no image is delivered and the package declares no asset dependency.  Editorial formatting only, in addition: an AMS \texttt{amsart} preamble that restates the article's own declarations and macros used by the retained lines, namely the environments \texttt{lem} on separate counters, together with the standard packages those lines need.  The retained environments print in this document as Lemma 1 in order of appearance, whereas the article prints the same items with the numbering of its own full document.  The complete native source of the article as distributed in the arXiv e-print for this exact version is bundled unchanged as \texttt{sources/173630/source0.tex}.
\begin{lem} \label{lem: n-1 2s}
    Let $\P$ be an $n$-dimensional Clements-Lindstr\"om poset, and $\Q$ a $(n-1)$-dimensional Clements-Lindstr\"om poset for $n > 2$. If $\P \diamond \Q$ is Macaulay, then $\P$ has exactly $n-1$ sides of length $2$ and $\Q$ has no sides of length $2$.
\end{lem}

\begin{lem} \label{lem: geq 5 not Mac}
Let $\P$ be a $n$-dimensional Clements-Lindstr\"om poset and $\Q$ a $(n - 1)$-dimensional Clements-Lindstr\"om poset for $n >2$. Then, $\P \diamond \Q$ is not Macaulay.
\end{lem}

\begin{proof}
Let $\P$ be the monomial poset of $k[x_1,\dots,x_n]/I$ where $I = (x_1^{d_1},\dots,x_n^{d_n})$ for $n \geq 5$. Also, let $\Q$ be the monomial poset of $k[y_1,\dots,y_{n-1}]/J$ where $J = (y_1^{c_1},\dots,y_{n-1}^{c_{n-1}})$ where for all $d_i \neq 1$, and $c_j \neq 1$, and $\P$ and $\Q$ are both of rank $s$. 
    
\emph{Case 1:} $n \geq 5$. 
Assume for the sake of contradiction that $\P \diamond \Q$ is Macaulay. Note that by \Cref{lem: n-1 2s}, $n-1$ of $d_1,\dots,d_n$ must be $2$, and also $c_j \neq 2$ for all $j$. Without loss of generality, let $d_1=\cdots=d_{n-1} = 2$, and $d_n \neq 2$. Assume that the maximal element of rank $1$ is $y_j$. Then, it follows as in the proof of \Cref{lem: geq 5 not Mac} that $\Seg_1 n-1 = \{y_1,\dots,y_{n-1}\}$. Note that $\Seg_1 n = \Seg_1 (n-1) \cup \{x_i\}$ for some $i$. Because of the definition of a diamond product, $\uSdw \{x_i\}$ must be distinct from $\uSdw \Seg_1 (n-1)$. Thus, we get that $|\uSdw \Seg_1 n| = n - 1 + \dots + 1 + n - 1$. Notably, if we take the set $A = \{x_1,\dots,x_n\}$ we get that $|\uSdw A| = n-1 + \dots + 1$. This contradicts that $y_j$ is the largest elements of rank $1$.
    
Now, assume that $x_i$ is the largest element of rank $1$.  Then we get that $\Seg_1 n = \{x_1, \dots, x_n\}$ by a similar argument as above. Now, by taking repeated upper shadows, we get that $\P_{[3]} = \Seg_3 t = \{x_1x_2x_3,\dots,x_n^3\}$, where $t$ is the number of elements of rank $3$ in $\P$. 
Using the hypothesis $n\geq 5$, we define a function $f: \Q_{[3]} \to \P_{[3]}$ by $f(y_jy_k) = f(y_j)f(y_k)$ for $j \neq k$, $f(y_j) = x_j$, $f(y_j^2) = x_jx_n$, and $f(y_j^3) = x_jx_n^2$ with the exception that $f(y_1^2y_2) = x_n^3$, $f(y_3^2y_j) = x_jx^2_n$, and $f(y_2^2y_4) = x_3x_n^2$. Note that $f$ is surjective, and because $f(y_1y_4^2) = f(y_1^2y_4)$, then $f$ is not injective. Thus, 
$|\Q_{[3]}| > |\P_{[3]}|$. Also, because $|\P_{[s-1]}| = n$, and $|\Q_{[s-1]}| = n-1$, there must exist some $k$ where $|\P_{[k]}| < |\Q_{[k]}|$, but $|\P_{[k+1]}| \geq |\Q_{[k+1]}|$. Note that $\P_{[k]} = \Seg_k |\P_{[k]}|$. Also,  $A= \Seg_k (|\P_{[k]}| + 1)$ must contain some element in $\Q$ of rank $k$. Because all elements of $\Q$ have a distinct upper shadow from the elements in $\P$, it follows that 
\begin{equation*}
|\uSdw A| > |\uSdw \P_{[k]}| =|\P_{[k+1]}|\geq  |\Q_{[k+1]}| \geq |\uSdw \Q_{[k]}|.
\end{equation*}
Thus, $|\uSdw \Seg_k |\Q_{[k]}|| \geq |\uSdw \Seg_k (|\P_{[k]}| + 1)| > |\uSdw \Q_{[k]}|$. This contradicts that $\P \diamond \Q$ is Macaulay.
      
\emph{Case 2:} $3 \leq n\leq 4$.

Note that the largest element of rank $1$ must be in $\P$. Without loss of generality let the largest element of rank $1$ be $x_1$. Assume that $n = 3$. Then, $\Seg_1 2 = \{x_1,x_2\}$. Assume that $c_1 = 3$, and $c_2 = m$. Because $\P$ and $\Q$ are of the same rank we have $3 + m = 2 + 2 + d_3 - 1$. Thus, $d_3 = m$. Now, by taking repeated upper shadows of $\Seg_1 2$, we get that $\Seg_{m-1} 3 = \{x_1x_2x_3^{m-3},x_1x_3^{m-2},x_2x_3^{m-2}\}$, and that $\uSdw \Seg_{m-1} 3 = \{x_1x_2x_3^{m-2},x_1x_3^{m-1},x_2x_3^{m-1}\}$. Also, if we take the set $A = \{y_1^2y_2^{m-3},y_1y_2^{m-2},y_2^{m-1}\}$, then $\uSdw A = \{y_1^2y_2^{m-2},y_1y_2^{m-1}\}$. This contradicts that $P \diamond Q$ is Macaulay. Thus, $c_1 > 3$. 
 
 Assume that $c_1 \geq 5$, and $c_1 < c_2$. Then we get that $\{y_1^4y_2,y_1^3y_2^2,y_1^2y_2^3,y_1y_2^4,y_2^5\} \subseteq \Q_{[5]}$, and  $\P_{[5]} = \{x_1x_2x_3^3,x_1x_3^4,x_2x_3^4,x_3^5\}$. Then, because $|\P_{[s-1]}| < |\Q_{[s-1]}|$ it follows that $\P \diamond \Q$ is not Macaulay. Thus, if $c_1 = 5$, then $c_2 = 5$. Then, note that $\Q_{[5]} = \{y_1^4y_2,y_1^3y_2^2,y_1^2y_2^3,y_1y_2^4\}$, and $\uSdw \Q_{[5]} = \{y_1^4y_2^2,y_1^3y_2^3,y_1^2y_2^4\}$. Moreover, because $\P$ and $\Q$ are both same rank, then $d_3 = 7$. Therefore, it follows that $\Seg_5 4 = \P_{[5]} = \{x_1x_2x_3^3,x_1x_3^4,x_2x_3^4,x_3^5\}$ and $\uSdw \P_{[5]} = \{x_1x_2x_3^4, x_1x_3^5, x_2x_3^5, x_3^6\}$, but this contradicts that $\P \diamond \Q$ is Macaulay. Thus, $c_1 = 4$. Also, let $c_2 = m$. Because $4 + m = 2 + 2 + d_3 - 1$ it follows that $d_3 = m+1$. Then
 \begin{eqnarray*}
 \P_{[m-1]} &=& \Seg_{m-1} 4 = \{x_1x_2x_3^{m-3},x_1x_3^{m-2},x_2x_3^{m-2},x_3^{m-1}\}\\
 \uSdw \P_{[m-1]} &=&\{x_1x_2x_3^{m-2},x_1x_3^{m-1},x_2x_3^{m-1},x_3^m\}\\
  \Q_{[m-1]} &=& \{y_1^3y_2^{m-4},y_1^2y_2^{m-3},y_1y_2^{m-2},y_2^{m-1}\}\\
  \uSdw \Q_{[m-1]} &=&\{y_1^3y_2^{m-3},y_1^2y_2^{m-2},y_1y_2^{m-1}\}.
  \end{eqnarray*}
This contradicts that $P \diamond Q$ is Macaulay. If $n = 3$, then $P \diamond Q$ is not Macaulay. Now assume that $n = 4$. Note that if $c_3 > 3$, then 
\begin{eqnarray*}
\Q_{[3]} &=&  \{y_1y_2y_3,y_1^2y_2,y_1^2y_3,y_1y_2^2,y_2^2y_3,y_1y_3^2,y_2y_3^2,y_3^3\}\\
\P_{[3]} &=&  \{x_1x_2x_3,x_1x_2x_4,x_1x_3x_4,x_2x_3x_4,x_1x_4^2,x_2x_4^2,x_3x_4^2\}.
\end{eqnarray*}
Then, because $|\Q_{[3]}| > |\P_{[3]}|$, and $|\Q_{[s-1]}| < |\P_{[s-1]}|$ means that $\P \diamond \Q$ is not Macaulay. Thus, $c_1 = c_2 = c_3 = 3$. Thus, because $3 + 3 + 3 = 2 + 2 + 2 + d_4 - 1$, then $d_4 = 2$ which contradicts \Cref{lem: n-1 2s}. Thus, if $n = 4$, then $\P \diamond \Q$ is not Macaulay. Therefore, for $n > 2$, $\P \diamond \Q$ is not Macaulay.
\end{proof}

\end{document}
